\documentclass[11pt]{article}

\usepackage[margin=1in]{geometry}
\usepackage{titlesec}
\usepackage{xcolor}
\usepackage{hyperref}
\usepackage{fontspec}
\usepackage{fancyhdr}
\usepackage{graphicx}
\usepackage{enumitem}
\usepackage{tcolorbox}
\usepackage{tikz}
\tcbuselibrary{skins,breakable}

\graphicspath{{manual_images/}}

\definecolor{brandgreen}{HTML}{059669}
\definecolor{brandteal}{HTML}{0d9488}
\definecolor{darktext}{HTML}{1f2937}
\definecolor{graytext}{HTML}{4b5563}
\definecolor{lightbg}{HTML}{f0fdf4}
\definecolor{bordergreen}{HTML}{a7f3d0}

% Noto Sans Bengali (bundled locally in fonts/) covers Bangla + basic
% Latin/numerals with real Regular and Bold weights, so it's used as the
% single main font rather than juggling a separate Latin face for the odd
% English UI term (Chat, Send, Edit) that appears inline with Bangla
% sentences. No dedicated italic file ships for this family, so the italic
% shape is a slanted synthesis of the regular weight rather than a true
% italic design.
\setmainfont{NotoSansBengali}[
  Path = fonts/,
  Extension = .ttf,
  UprightFont = *-Regular,
  BoldFont = *-Bold,
  ItalicFont = *-Regular,
  ItalicFeatures = {FakeSlant=0.2},
]

\titleformat{\section}
  {\normalfont\Large\bfseries\color{brandgreen}}
  {\thesection}{0.6em}{}
\titleformat{\subsection}
  {\normalfont\large\bfseries\color{darktext}}
  {\thesubsection}{0.5em}{}

\hypersetup{
  colorlinks=true,
  linkcolor=brandgreen,
  urlcolor=brandteal,
  pdftitle={KenaBecha JU -- ব্যবহারকারী নির্দেশিকা},
  pdfauthor={KenaBecha JU}
}

\pagestyle{fancy}
\setlength{\headheight}{14pt}
\fancyhf{}
\renewcommand{\headrulewidth}{0.4pt}
\fancyhead[L]{\small\color{graytext} KenaBecha JU}
\fancyhead[R]{\small\color{graytext} ব্যবহারকারী নির্দেশিকা}
\fancyfoot[C]{\small\color{graytext} \thepage}

\setlist[itemize]{leftmargin=1.4em, itemsep=3pt, topsep=3pt}
\setlist[enumerate]{leftmargin=1.6em, itemsep=4pt, topsep=4pt}

% A screenshot in a soft-bordered card, with a small gray caption below --
% the same visual language (rounded, soft green border) as the app's own
% cards, so the manual doesn't feel like a separate, colder document.
\newtcolorbox{shotbox}{
  colback=white,
  colframe=bordergreen,
  boxrule=1pt,
  arc=3mm,
  boxsep=2pt,
  left=2pt, right=2pt, top=2pt, bottom=2pt,
  breakable
}
\newcommand{\screenshot}[2]{
  \begin{center}
    \begin{shotbox}
      \includegraphics[width=0.92\linewidth]{#1}
    \end{shotbox}

    \vspace{2pt}
    {\small\color{graytext}\itshape #2}
  \end{center}
  \vspace{4pt}
}
% A small colored callout box for tips/notes, matching the app's own
% amber/blue notice-banner convention rather than a plain LaTeX quote block.
\newtcolorbox{tipbox}{
  colback=lightbg, colframe=bordergreen, boxrule=0.8pt, arc=2mm,
  left=8pt, right=8pt, top=6pt, bottom=6pt, breakable
}

\begin{document}

\begin{titlepage}
  \centering
  \vspace*{2.2cm}
  {\Huge\bfseries\color{brandgreen} KenaBecha JU}\\[0.4cm]
  {\LARGE\color{darktext} ব্যবহারকারী নির্দেশিকা}\\[0.3cm]
  {\large\color{graytext} জাহাঙ্গীরনগর বিশ্ববিদ্যালয়ের ক্যাম্পাস মার্কেটপ্লেস}\\[1.6cm]

  {\normalsize\color{graytext}
  এই বইটি দুইটি অংশে ভাগ করা --- প্রথম অংশ তাদের জন্য যারা \textbf{কিনতে} চান,
  দ্বিতীয় অংশ তাদের জন্য যারা নিজের \textbf{দোকান চালিয়ে বিক্রি} করতে চান।
  প্রতিটি ধাপ আসল স্ক্রিনশট দিয়ে দেখানো হয়েছে, যাতে বইয়ের সাথে সাথে
  আপনি নিজেও করে দেখাতে পারেন।}

  \vfill
  {\small\color{graytext} \today}
\end{titlepage}

\tableofcontents
\newpage

\section*{ভূমিকা}
\addcontentsline{toc}{section}{ভূমিকা}

কেনাবেচা জেইউ (KenaBecha JU) জাহাঙ্গীরনগর বিশ্ববিদ্যালয়ের শিক্ষার্থীদের
জন্য তৈরি একটি নিজস্ব মার্কেটপ্লেস। এখানে আপনি ক্যাম্পাসের অন্য
শিক্ষার্থীদের কাছ থেকে বই, ইলেকট্রনিক্স, হোস্টেলের জিনিসপত্র --- প্রায়
সবকিছু কিনতে ও বিক্রি করতে পারবেন। শুধু তা-ই না, চাইলে নিজের একটা
``দোকান'' খুলে নিয়মিত ব্যবসাও চালাতে পারবেন, ঠিক যেভাবে ক্যাম্পাসের
বাইরে কেউ একটা দোকান চালায়।

প্রতিটি বিক্রেতা জাবি-ভেরিফাইড, অর্থাৎ যার সাথে আপনি কথা বলছেন তিনি
সত্যিকারের একজন শিক্ষার্থী --- অচেনা কেউ নন। টাকা-পয়সার লেনদেন
অ্যাপের ভেতরে হয় না; আপনি নিজেরা কথা বলে ক্যাম্পাসের মধ্যেই দেখা করে
কেনাবেচা সম্পন্ন করবেন।

\begin{tipbox}
\textbf{একটা কথা মনে রাখবেন:} এই বইয়ের সব স্ক্রিনশট বাংলা ভাষায় নেওয়া।
অ্যাপের উপরের ডানদিকে একটা \textbf{ভাষা বাটন} আছে (গ্লোব আইকন) ---
সেখানে চাপলে যেকোনো সময় ইংরেজি বা বাংলার মধ্যে পাল্টাতে পারবেন।
\end{tipbox}

\screenshot{01_home}{কেনাবেচা জেইউ-এর হোমপেজ --- বাংলা ভাষায়}

\newpage
\part*{অংশ ১ --- ক্রেতাদের জন্য}
\addcontentsline{toc}{section}{অংশ ১ --- ক্রেতাদের জন্য}

এই অংশে দেখানো হয়েছে কীভাবে অ্যাকাউন্ট খুলবেন, পণ্য খুঁজবেন, বিক্রেতার
সাথে কথা বলবেন, এবং কেনার পর রিভিউ দেবেন।

\section{অ্যাকাউন্ট তৈরি করা}

প্রথমেই একটা অ্যাকাউন্ট লাগবে। হোমপেজের ওপরে ডানদিকে \textbf{সাইন আপ}
বাটনে চাপুন। দুইটা উপায় আছে:

\begin{itemize}
  \item শুধু কেনাকাটা করতে চাইলে --- \textbf{Google দিয়ে সাইন ইন} করে
  এক ক্লিকেই ঢুকে পড়তে পারবেন।
  \item পরে যদি নিজেও কিছু বিক্রি করতে চান, তাহলে প্রথম থেকেই
  \textbf{``সম্পূর্ণ জাবি তথ্য দিয়ে সাইন আপ''} বেছে নেওয়া ভালো ---
  কারণ বিক্রি করার জন্য আপনার হল, বিভাগ, সেশন ইত্যাদি তথ্য লাগবেই।
\end{itemize}

ফর্মে আপনার নাম, ইমেইল, পাসওয়ার্ড, ফোন নম্বর, স্টুডেন্ট আইডি,
রেজিস্ট্রেশন নম্বর, হল, বিভাগ ও সেশন দিতে হবে। এই তথ্যগুলো যাচাই করেই
বোঝা যায় যে আপনি সত্যিকারের একজন জাবি শিক্ষার্থী।

\screenshot{03_signup_form}{সাইন আপ ফর্ম --- নাম, ইমেইল, ফোন, হল, বিভাগ ও সেশন পূরণ করা হচ্ছে}

\section{ইমেইল যাচাই (OTP)}

ফর্ম জমা দেওয়ার পর আপনার ইমেইলে একটা ৬-সংখ্যার কোড (OTP) পাঠানো হবে।
এই পাতায় সেই ৬টা সংখ্যা একটার পর একটা লিখলেই আপনার ইমেইল যাচাই হয়ে
যাবে। কোড না পেলে কিছুক্ষণ পর ``আবার পাঠান'' বাটনে চাপতে পারবেন।

\screenshot{04_verify_email}{কোড দেওয়ার পাতা}
\screenshot{05_otp_filled}{৬টি সংখ্যা দেওয়া হয়েছে}

যাচাই হয়ে গেলে আপনাকে লগইন পাতায় পাঠানো হবে --- সেখানে ইমেইল ও
পাসওয়ার্ড দিয়ে লগইন করলেই আপনি ভেতরে ঢুকে যাবেন।

\screenshot{06_login}{লগইন পাতা --- ইমেইল যাচাই সফল হয়েছে, এখন লগইন করতে হবে}

\section{পণ্য খোঁজা ও ব্রাউজ করা}

লগইন করার পর হোমপেজের ওপরে \textbf{পণ্যসমূহ} মেনুতে চাপলে সব পণ্যের
তালিকা দেখতে পাবেন। প্রতিটা কার্ডে ছবি, নাম ও দাম দেখানো থাকে।

\screenshot{16_browse_listings}{সব পণ্যের তালিকা --- প্রতিটা কার্ডে ছবি ও দাম দেখা যাচ্ছে}

নির্দিষ্ট কিছু খুঁজতে চাইলে ওপরের সার্চ বক্সে লিখে খুঁজুন --- যেমন
নিচের উদাহরণে ``বই'' লিখে সার্চ করা হয়েছে।

\screenshot{17_search_results}{``বই'' লিখে সার্চ করার ফলাফল}

\begin{tipbox}
পণ্যের পাশাপাশি \textbf{দোকানও} খুঁজতে পারবেন --- ওপরের মেনুতে
\textbf{দোকান ব্রাউজ করুন}-এ চাপলে ক্যাম্পাসের সব দোকানের তালিকা
দেখা যাবে।
\end{tipbox}

\section{পণ্যের বিস্তারিত দেখা}

কোনো পণ্যে চাপলে তার নিজস্ব পাতা খুলবে --- সেখানে ছবি (একাধিক থাকলে
সোয়াইপ করে দেখা যায়), দাম, অবস্থা (নতুন/পুরাতন), পণ্য কীভাবে হাতে
পাবেন (পিকআপ/ডেলিভারি/উভয়) এবং পিকআপের ঠিকানা --- সবকিছু বিস্তারিত
দেখা যায়।

\screenshot{18_listing_detail_buyer}{একটা পণ্যের বিস্তারিত পাতা --- দাম, ঠিকানা ও ``বিক্রেতার সাথে চ্যাট করুন'' বাটন}

\section{বিক্রেতার সাথে চ্যাট করা}

পণ্য পছন্দ হলে \textbf{Chat with seller} বাটনে চাপুন --- সাথে সাথে
বিক্রেতার সাথে একটা চ্যাট থ্রেড খুলে যাবে, যেখানে সরাসরি লেখালেখি করে
সময়-জায়গা ঠিক করতে পারবেন।

\screenshot{19_chat_started}{চ্যাট শুরু হয়েছে --- পণ্যটা সংযুক্ত অবস্থায় দেখা যাচ্ছে}

নিচের বক্সে বার্তা লিখে \textbf{Send} বাটনে চাপুন বা কি-বোর্ডে Enter
চাপুন। বার্তা তখনই পাঠানো হয়ে যায় ও বিক্রেতা রিয়েল-টাইমে দেখতে পান।

\screenshot{20_chat_message_sent}{বার্তা পাঠানো হয়েছে --- বিক্রেতাকে জিজ্ঞেস করা হচ্ছে বইগুলো এখনও আছে কিনা}

\section{বিক্রেতার উত্তর ও নোটিফিকেশন}

বিক্রেতা উত্তর দিলে আপনি সাথে সাথে একটা নোটিফিকেশন পাবেন --- ওপরের
ঘণ্টা আইকনে একটা লাল সংখ্যা দেখা যাবে। সেখানে চাপলে সাম্প্রতিক সব
নোটিফিকেশনের তালিকা দেখা যায়।

\screenshot{27_buyer_notification_bell}{নোটিফিকেশন প্যানেল --- বিক্রেতার নতুন বার্তার নোটিফিকেশন}

\section{কেনার পর রিভিউ দেওয়া}

জিনিসটা হাতে পাওয়ার পর, বিক্রেতা যখন পণ্যটা ``বিক্রিত'' হিসেবে
চিহ্নিত করে দেবেন, তখন সেই পণ্যের পাতাতেই একটা \textbf{রিভিউ দেওয়ার}
ফর্ম দেখা যাবে --- ১ থেকে ৫ তারা দিয়ে রেটিং ও ইচ্ছা করলে একটা লেখা
মন্তব্যও যোগ করতে পারবেন।

\screenshot{28_rate_filled}{রিভিউ ফর্ম --- ৫ তারা ও একটা মন্তব্য লেখা হয়েছে}

\begin{tipbox}
রিভিউ দেওয়ার অপশনটা তখনই দেখা যায় যখন সত্যিই আপনি সেই বিক্রেতার
সাথে বার্তা চালাচালি করেছেন এবং পণ্যটা বিক্রিত হিসেবে চিহ্নিত হয়েছে
--- এলোমেলোভাবে যেকোনো পণ্যে রিভিউ দেওয়া যায় না, এতে রিভিউগুলো
সবসময় আসল অভিজ্ঞতার উপর ভিত্তি করেই হয়।
\end{tipbox}

আপনার দেওয়া রিভিউ সাথে সাথে সেই দোকানের পাবলিক পাতায় দেখা যেতে শুরু
করে, যাতে অন্য ক্রেতারাও বুঝতে পারেন এই দোকান কেমন।

\screenshot{29_shop_with_review}{দোকানের পাতায় নতুন রিভিউ ও ৫.০ রেটিং দেখা যাচ্ছে}

\section{এআই শপিং সহকারী}

কিছু খুঁজে না পেলে বা প্রশ্ন থাকলে, পাতার নিচে-ডানদিকে থাকা সবুজ গোল
বাটনে চাপুন --- এটাই কেনাবেচা জেইউ-এর নিজস্ব এআই শপিং সহকারী। একে
সরাসরি বাংলায় লিখে জিজ্ঞেস করা যায়, যেমন কোনো দোকানের নাম বা কোনো
পণ্যের ধরন।

\screenshot{30_ai_assistant_reply}{এআই সহকারীকে ``ক্যাম্পাস কর্নার'' লিখে জিজ্ঞেস করা হয়েছে, এবং সে আসল দোকানটা খুঁজে দেখাচ্ছে}

সহকারী শুধু ওয়েবসাইটে সত্যিই যা আছে তা-ই দেখায় --- কোনো পণ্য বা
দোকান বানিয়ে বলে না। তাই এখানে যা দেখানো হয়, তা সবসময় সত্যিকারের।

\section{নিজের প্রোফাইল}

ওপরের ডানদিকে নিজের ছবি/আইকনে চাপলে একটা মেনু খুলবে, সেখান থেকে
\textbf{প্রোফাইল} নির্বাচন করলে আপনার নিজের তথ্য, রেটিং ও কেনাবেচার
ইতিহাস দেখা যাবে।

\screenshot{31_account_menu}{অ্যাকাউন্ট মেনু}
\screenshot{32_profile_page}{নিজের প্রোফাইল পাতা}

\newpage
\part*{অংশ ২ --- উদ্যোক্তাদের জন্য}
\addcontentsline{toc}{section}{অংশ ২ --- উদ্যোক্তাদের জন্য}

এই অংশে দেখানো হয়েছে কীভাবে নিজের দোকান খুলবেন, পণ্য যোগ করবেন,
পোস্ট দিয়ে প্রচার করবেন এবং ক্রেতাদের সাথে লেনদেন সামলাবেন। অংশ ১-এ
দেখানো অ্যাকাউন্ট তৈরির প্রক্রিয়াটা এখানেও একই --- শুধু পার্থক্য হলো
আপনাকে এবার নিজের দোকান খুলতে হবে।

\section{নিজের দোকান খোলা}

লগইন করার পর ওপরের মেনু থেকে \textbf{আমার দোকান}-এ যান। কোনো দোকান
না থাকলে একটা খালি পাতা দেখা যাবে ---সেখান থেকে \textbf{নতুন দোকান}
বাটনে চাপুন।

\screenshot{07_open_shop_empty}{আমার দোকান পাতা --- এখনও কোনো দোকান খোলা হয়নি}

দোকানের নাম, ধরন (যেমন Books \& Study) এবং সংক্ষিপ্ত বিবরণ লিখুন।
চাইলে একটা লোগো ছবিও যোগ করতে পারেন। এই তথ্যগুলোই আপনার দোকানের
পাবলিক পাতায় দেখা যাবে।

\screenshot{08_new_shop_form}{নতুন দোকান তৈরির ফর্ম --- নাম, ধরন ও বিবরণ পূরণ করা হয়েছে}

দোকান তৈরি হয়ে গেলে সেটা আপনার ড্যাশবোর্ডে দেখা যাবে, সাথে
\textbf{পণ্য যোগ করুন} ও \textbf{Posts} নামের দুইটা বাটন।

\screenshot{09_shop_created}{দোকান তৈরি হয়ে গেছে --- ``ক্যাম্পাস কর্নার''}

\section{পণ্য যোগ করা}

\textbf{পণ্য যোগ করুন} বাটনে চাপলে একটা বিস্তারিত ফর্ম খুলবে। এখানে
লিখতে হবে:

\begin{itemize}
  \item পণ্যের নাম ও বিবরণ
  \item ক্যাটাগরি (যেমন Books \& Study)
  \item দাম ও একক (যেমন প্রতি সেট, প্রতি কেজি)
  \item এক বা একাধিক ছবি
  \item ক্রেতা পণ্যটা কীভাবে পাবেন --- \textbf{পিকআপ}, \textbf{ডেলিভারি}, নাকি
  \textbf{উভয়ই} (ক্রেতা নিজে বেছে নেবেন)
  \item পিকআপের ঠিকানা (উভয়ই বা শুধু পিকআপ বেছে নিলে)
\end{itemize}

\screenshot{11_add_listing_filled}{পণ্য যোগ করার ফর্ম --- নাম, দাম, হস্তান্তর পদ্ধতি ও পিকআপ ঠিকানা পূরণ করা হয়েছে, দুইটা ছবি যোগ করা হয়েছে}

জমা দেওয়ার সাথে সাথে পণ্যটা লাইভ হয়ে যায় --- ক্রেতারা এখন থেকেই
সেটা খুঁজে পাবেন ও কিনতে পারবেন।

\screenshot{12_listing_owner_view}{পণ্যের পাতা, মালিকের দৃষ্টিতে --- এখানে Edit ও Delete অপশন দেখা যায়}

\begin{tipbox}
নিজের দোকানের পণ্যে আপনি \textbf{Edit} ও \textbf{Delete} বাটন
দেখবেন --- কিন্তু কোনো ক্রেতা এই একই পাতায় গেলে তার বদলে
\textbf{Chat with seller} বাটন দেখবেন। পাতাটা একই, শুধু আপনি কে
তার ওপর নির্ভর করে দেখানো অপশন পাল্টে যায়।
\end{tipbox}

\section{পণ্য নিয়ে পোস্ট দেওয়া}

শুধু পণ্য যোগ করাই যথেষ্ট নয় --- মাঝে মাঝে সেই পণ্য নিয়ে একটা পোস্ট
দিলে বেশি মানুষের নজরে পড়ে। দোকান ড্যাশবোর্ডে \textbf{Posts}-এ গিয়ে
\textbf{New post}-এ চাপুন।

\screenshot{14_new_post_empty}{পোস্ট পাতা --- এখনও কোনো পোস্ট নেই}

শিরোনাম, লেখা এবং ছবি দিন --- চাইলে আপনার দোকানের কোনো পণ্যও পোস্টের
সাথে যুক্ত করতে পারেন, যাতে পাঠক সরাসরি সেই পণ্যের পাতায় যেতে পারেন।

\screenshot{15_post_filled}{নতুন পোস্ট --- শিরোনাম, লেখা ও একটা ছবি যোগ করা হয়েছে}

\begin{tipbox}
প্রতিটা পোস্ট প্রকাশ হওয়ার আগে অ্যাডমিনের অনুমোদন লাগে --- এটা
স্প্যাম ঠেকানোর জন্য। অনুমোদন পেলেই পোস্টটা সবার ফিডে দেখা যাবে,
এবং যারা আপনার দোকান ফলো করেন তারা একটা নোটিফিকেশনও পাবেন।
\end{tipbox}

\section{দোকানের পাবলিক পাতা}

আপনার দোকানের নিজস্ব একটা পাবলিক পাতা আছে, যেখানে যে কেউ গিয়ে আপনার
সব পণ্য, পোস্ট ও রিভিউ একসাথে দেখতে পারেন।

\screenshot{13_shop_storefront}{দোকানের পাবলিক পাতা --- পণ্যের সংখ্যা, বিক্রির সংখ্যা ও রেটিং একসাথে দেখা যাচ্ছে}

\section{ক্রেতার বার্তার উত্তর দেওয়া}

কোনো ক্রেতা আপনাকে বার্তা পাঠালে ওপরের ঘণ্টা আইকনে নোটিফিকেশন আসবে।

\screenshot{21_seller_notification_bell}{নতুন বার্তার নোটিফিকেশন}

\textbf{ইনবক্স}-এ গেলে সব কথোপকথনের তালিকা দেখা যাবে।

\screenshot{22_seller_inbox_list}{ইনবক্স --- ক্রেতার নতুন বার্তা দেখা যাচ্ছে}

কথোপকথনে ঢুকে সরাসরি উত্তর লিখে পাঠাতে পারবেন --- ঠিক যেভাবে ক্রেতা
আপনাকে লিখেছিলেন।

\screenshot{23_seller_chat_reply}{ক্রেতাকে উত্তর দেওয়া হয়েছে --- দেখা করার সময় ঠিক করা হচ্ছে}

\section{পণ্য বিক্রিত হিসেবে চিহ্নিত করা}

জিনিসটা সত্যিই বিক্রি হয়ে গেলে, সেই পণ্যের \textbf{Edit} পাতায় গিয়ে
নিচের দিকে ``স্ট্যাটাস'' অংশে \textbf{বিক্রিত হিসেবে চিহ্নিত করুন}
বাটনে চাপুন।

\screenshot{24_listing_edit_status}{পণ্য সম্পাদনার পাতা --- স্ট্যাটাস পরিবর্তনের বাটনগুলো নিচের দিকে}

এতে পণ্যটা তালিকা থেকে সরে যায় না, বরং ``বিক্রিত'' ব্যাজ নিয়ে
দেখানো হয় --- এবং এখনই ক্রেতার জন্য রিভিউ দেওয়ার অপশনটা চালু হয়ে
যায়।

\screenshot{26_listing_sold_public}{পণ্যের পাতায় এখন ``বিক্রিত'' ব্যাজ দেখা যাচ্ছে}

\begin{tipbox}
নিজের একক (personal) বিক্রির পণ্যের জন্য এই বাটনটা সরাসরি পণ্যের
পাতাতেই পাওয়া যায়। কিন্তু দোকানের পণ্যের জন্য এটা Edit পাতায় থাকে
--- কারণ দোকানের পণ্যে আরও বেশি স্ট্যাটাস অপশন থাকে, যেমন
\textit{স্টক শেষ}, \textit{বিরতি} ও \textit{আবার চালু করুন}।
\end{tipbox}

\section{রিভিউ পাওয়া}

ক্রেতা যখন রিভিউ দেন, সেটা সাথে সাথে আপনার দোকানের পাবলিক পাতায় ও
সামগ্রিক রেটিং-এ যোগ হয়ে যায় --- নতুন কোনো ক্রেতা আপনার দোকানে ঢুকলে
প্রথমেই এই রেটিং দেখতে পাবেন।

\screenshot{29_shop_with_review}{দোকানের পাতায় প্রথম রিভিউ যোগ হয়েছে}

\newpage
\section*{শেষ কথা}
\addcontentsline{toc}{section}{শেষ কথা}

এই বইয়ে দেখানো প্রতিটা ধাপ --- সাইন আপ থেকে শুরু করে দোকান খোলা,
পণ্য বিক্রি, চ্যাট, রিভিউ এবং এআই সহকারী পর্যন্ত --- আসল অ্যাপে সরাসরি
করে দেখানো হয়েছে, কোনো কল্পিত বর্ণনা নয়। তাই এখানে যা দেখলেন,
বাস্তবে ব্যবহার করার সময়ও ঠিক একইরকম অভিজ্ঞতা পাবেন।

কোনো সমস্যা হলে বা প্রশ্ন থাকলে সাইটের নিচে থাকা \textbf{যোগাযোগ}
তথ্য থেকে অ্যাডমিনের সাথে যোগাযোগ করতে পারেন।

\vspace{1cm}
\begin{center}
{\Large\color{brandgreen}\bfseries শুভকামনা, ক্যাম্পাসেই কিনুন, বিক্রি করুন!}
\end{center}

\end{document}
